\section{Legal Landscape}
\label{sec:legal}

\subsection{\emph{Reynolds v. Sims} and the Equal Population Mandate}

The modern doctrine begins with \emph{Reynolds v. Sims}, 377 U.S. 533 (1964), where Chief Justice Warren declared that ``legislators represent people, not trees or acres.'' The Court held that state legislative districts must contain substantially equal populations, establishing the constitutional foundation for the one-person, one-vote principle. Critically, the Court did not define which population. The decennial Census total-population count was available and administratively convenient; it became the default by practice rather than by mandate.

For congressional districts specifically, \emph{Wesberry v. Sanders}, 376 U.S. 1 (1964), decided two months before \emph{Reynolds}, grounded the equal-population requirement in Article I, Section 2 of the Constitution, which requires that Representatives be chosen ``by the People.'' Justice Black's majority opinion read ``the People'' to require numerical equality in the apportionment of House seats. Again, no definition of which population unit was specified.

\subsection{Intermediate Developments}

For four decades after \emph{Reynolds} and \emph{Wesberry}, the question of which population to equalize was largely dormant. States uniformly used Census total population; no party litigated an alternative. Two developments changed this landscape.

First, the growth of non-citizen immigrant populations, particularly in California, Texas, and New York, created demographic conditions where total population and CVAP diverged substantially. By 2010, several California congressional districts had non-citizen adult populations exceeding 20\% of the total, meaning that the districts' ``people'' in the \emph{Reynolds} sense were heavily weighted toward non-voters.

Second, political mobilization by conservative litigants and scholars who argued that non-citizen counting artificially inflated the congressional representation of high-immigration states and districts, predominantly Democratic-leaning. The Public Interest Legal Foundation and related groups began supporting litigants who challenged total-population apportionment as constitutionally impermissible.

\subsection{\emph{Evenwel v. Abbott} (2016)}

The Supreme Court's unanimous decision in \emph{Evenwel v. Abbott}, 578 U.S. 54 (2016), is the most important modern precedent on balance-metric choice. Plaintiffs in Texas argued that equalizing total population (rather than eligible voters) diluted the weight of their votes by packing their district with non-voters, violating the equal-protection guarantee of \emph{Reynolds}.

Justice Ginsburg, writing for a unanimous Court, held that states \emph{may} use total population as their apportionment base and that doing so does not violate the Equal Protection Clause. The decision rested on three grounds: (1) historical practice dating to the founding; (2) the whole-person representation theory---elected officials serve all constituents, not just voters; and (3) the administrative convenience of the decennial Census.

Critically, the Court explicitly declined to decide whether states \emph{must} use total population or whether they \emph{may} use an alternative metric such as CVAP. Justice Thomas concurred in the judgment but wrote separately to argue that states should be free to use any permissible metric. Justice Alito also concurred, agreeing the case was not the right vehicle to resolve the ultimate question.

The result is a permissive framework: total population is constitutionally permissible; CVAP is probably permissible (at least for state legislative districts); no metric has been ruled impermissible by the Supreme Court except to the extent that a metric produces individual-district deviations exceeding acceptable tolerances.

\subsection{State-Level Legal Landscape}

No state currently uses CVAP as its primary congressional redistricting metric; total population is universal for congressional districts. However, several states have experimented with alternative metrics for state legislative redistricting:

\begin{itemize}
    \item \textbf{California}: The Citizens Redistricting Commission has discussed CVAP-adjusted plans but ultimately uses total population.
    \item \textbf{Utah}: In 2021, the legislature passed H.B. 131 directing the Independent Redistricting Commission to report CVAP-equivalent maps alongside total-population maps for informational purposes.
    \item \textbf{Missouri}: Courts have considered registered-voter arguments in state legislative contexts.
\end{itemize}

The post-\emph{Evenwel} question is whether a future Court, with different membership, would affirmatively hold that states \emph{must} use CVAP or that total-population apportionment violates one-person-one-vote when it produces large CVAP disparities. Legal scholars are divided~\citep{scharff2016evenwel, ban2016counting}, but no current circuit split suggests imminent Supreme Court re-examination.

\subsection{Federal Apportionment: A Different Standard}

For purposes of \emph{apportioning} House seats among states, the Apportionment Clause of Article I requires using the ``whole number of persons in each state.'' The Census Bureau counts all residents including non-citizens for apportionment purposes, a practice affirmed in \emph{Dep't of Commerce v. New York}, 588 U.S. 752 (2019), which struck down the Trump administration's attempt to add a citizenship question to the 2020 Census (finding it was added for pretextual reasons, not invalidating citizenship-question authority as such). The district-level balance-metric question is therefore legally distinct from the apportionment-among-states question.
